\chapter{特征值与相似：寻找合适的基}\label{ch:eigen}
\chaptermap{特征值与谱}{相似与对角化}{实对称谱定理}{Jordan初步}
\begin{examtip}
\textbf{考研核心：}特征多项式、特征向量、迹与行列式、相似不变量、对角化及实对称矩阵的正交对角化。
\textbf{竞赛视角：}从特征空间维数研究结构，用矩阵多项式处理幂；重根不足以判断对角化，必须核算几何重数。Jordan内容标为选学提高。
\end{examtip}
\section{特征多项式与谱的基本关系}
\subsection{代数重数与几何重数}
\begin{definition}{特征值、特征空间与重数}
设$A\in\F^{n\times n}$。若存在$\lambda\in\F$及$\vx\in\F^n\setminus\{0\}$使$A\vx=\lambda\vx$，则$\lambda$称特征值，$\vx$称属于$\lambda$的特征向量。对一个特征值$\lambda$，特征空间$E_\lambda=\ker(A-\lambda I)$由其全部特征向量和零向量组成；零向量本身不是特征向量。$\chi_A(t)=\det(tI-A)$中$\lambda$作为根的重数称代数重数$a_\lambda$，$g_\lambda=\dim E_\lambda=n-\rank(A-\lambda I)$称几何重数。
\end{definition}
\begin{theorem}{谱恒等式}
在复数域按代数重数列全部特征值$\lambda_i$，有
\begin{equation}\sum_i\lambda_i=\tr A,\quad\prod_i\lambda_i=\det A,\quad
1\le g_\lambda\le a_\lambda.\label{eq:spectrum}
\end{equation}
若$A\vx=\lambda\vx$，则$p(A)\vx=p(\lambda)\vx$；$A$可逆时$A\inv$在该向量上的特征值为$\frac{1}{\lambda}$。
\end{theorem}
\begin{proof}
\label{der:5-1}
在$\det(tI-A)$中，要得到$t^{n-1}$项，必须有$n-1$个对角位置选择$t$，剩下的对角位置选择$-a_{ii}$；任何非恒等排列至少移动两个位置，不能给出此次数。因此该系数为$-\sum_i a_{ii}=-\tr A$。常数项为$\det(-A)=(-1)^n\det A$。与
\[\chi_A(t)=\prod_i(t-\lambda_i)=t^n-\left(\sum_i\lambda_i\right)t^{n-1}+\cdots+(-1)^n\prod_i\lambda_i\]
比较，得到迹与行列式关系。

特征值意味着$\ker(A-\lambda I)$含非零向量，故$g_\lambda\ge1$。取其$g_\lambda$个基向量扩充为全空间基，在这组基下$A$形如$\begin{pmatrix}\lambda I_g&B\\0&C\end{pmatrix}$，于是$\chi_A(t)=(t-\lambda)^g\det(tI-C)$，故代数重数至少$g$。

对特征向量，归纳有$A^k\vx=\lambda^k\vx$：$k=0$为恒等；若$k$成立，左乘$A$就得到$k+1$。线性组合即$p(A)\vx=p(\lambda)\vx$。若$A$可逆，$\lambda\ne0$，对$A\vx=\lambda\vx$左乘逆再除以$\lambda$，得$A\inv\vx=\lambda^{-1}\vx$。

为确认谱映射给出\emph{全部}特征值，可在复数域归纳上三角化：先取一个特征向量作首基向量，得到上述$g=1$的块形式；再在右下块递归，最终$P\inv AP=T$为上三角阵，对角元是全部特征值。$p(T)$仍上三角且对角元为$p(\lambda_i)$，因为三角矩阵乘积的对角元逐项相乘。又$p(A)=Pp(T)P\inv$，所以$p(A)$的全部谱值正是这些$p(\lambda_i)$，含重数。这同时证明$\tr(A^k)=\sum_i\lambda_i^k$，不要求可对角化。

\end{proof}
\begin{method}{特征值先用结构筛选}
三角矩阵直接读对角元；行和相同可试全1向量；低秩矩阵先找核。算出根后用迹、行列式核验。易错：对每个特征值必须解齐次方程，不能将任意非零向量配给该根。
\end{method}
\worked{三角矩阵的特征向量}{基础例题}{\star\star}{5-1}
{求$A=\begin{pmatrix}2&1\\0&3\end{pmatrix}$的特征值与特征空间。}
{特征值直接读，特征向量另解方程。}
{\textbf{方法一：}$\chi_A(t)=\det\begin{pmatrix}t-2&-1\\0&t-3\end{pmatrix}=(t-2)(t-3)$，根为2、3。根2时$(A-2I)(x,y)\T=(y,y)\T=0$，得$y=0,x$任意，所以$E_2=\Span(e_1)$；根3时$(A-3I)(x,y)\T=(-x+y,0)\T=0$，得$y=x$，所以$E_3=\Span((1,1)\T)$。
\textbf{方法二：}计算$Ae_1=(2,0)\T=2e_1$，$A(1,1)\T=(3,3)\T=3(1,1)\T$。两向量组成矩阵行列式1，独立且已有二维特征基，所以没有遗漏的特征方向或其他根。}
{观察可省行列式，方程法更系统。}
{特征空间写成所有倍数，特征向量需排除零倍数。}
\worked{迹与行列式的核验}{考研风格}{\star}{5-2}
{三阶矩阵特征值为$1,2,\lambda$且迹为7，求行列式与$\tr(A^2)$。}
{按代数重数计入全部根。}
{由迹关系$1+2+\lambda=7$，得$\lambda=4$。按代数重数取乘积，$\det A=1\cdot2\cdot4=8$。上三角化后$A$对角为$1,2,4$，平方的对角逐项平方，因此$\tr(A^2)=1^2+2^2+4^2=1+4+16=21$。这一推导不需要先假设对角化；本题三根互异虽确可对角化，但一般迹幂公式也成立。}
{迹可检查根的和，但一般不能唯一确定全部根。}
{$\tr(A^2)$通常不等于$(\tr A)^2$。}
\worked{多项式的特征值}{考研风格}{\star\star}{5-3}
{若$A$的特征值为$1,2,3$，求$A^2-3A+2I$的行列式。}
{把多项式作用于每个特征值。}
{\textbf{方法一：}令$p(t)=t^2-3t+2=(t-1)(t-2)$。谱映射给$p(A)$的根依次为$p(1)=0,p(2)=0,p(3)=9-9+2=2$，因此$\det p(A)=0\cdot0\cdot2=0$。
\textbf{方法二：}$p(A)=(A-I)(A-2I)$，因为$I$与$A$可交换。$A$含特征值1和2，分别使$A-I,A-2I$奇异，两个行列式为零；乘积公式给$\det p(A)=\det(A-I)\det(A-2I)=0$。第二法只需识别一个零因子就足够。}
{只求行列式时无需特征向量。}
{$p(A)$可能比$A$有更大的特征空间，不能据新谱反推旧结构。}
\worked{秩一矩阵的谱}{竞赛风格}{\star\star\star}{5-4}
{非零列向量$\vu,\vv\in\R^n$，研究$K=\vu\vv\T$的特征值及其可对角化条件。}
{唯一可能非零特征值是$\vv\T\vu=c$。}
{先有$Kx=u(v\T x)$，所以像包含$\Span(u)$，且$u,v$非零保证$K$秩1。记$c=v\T u$，则$Ku=cu$、$K^2=cK$。
\textbf{方法一：}若$c\ne0$，$u\notin\ker v\T$；两空间维数为1和$n-1$，交为零，构成直和。在$\Span(u)$上特征值$c$，在$\ker v\T$上为0，拼基得到可对角化。若$c=0$，$K^2=0$且$K\ne0$，任一特征值满足$\lambda^2=0$，全为0；若可对角化就相似于零阵，迫使$K=0$，矛盾。
\textbf{方法二：}当$c\ne0$，令$P=\frac{K}{c}$，则$P^2=P$，故空间分解成$\im P\oplus\ker P$，其上$K$分别乘$c,0$。等价地，零化多项式$t(t-c)$互异一次因子保证可对角化。$c=0$时最低零化多项式为$t^2$，有重因子，得到同样否定结论。}
{空间法直接构造特征基；多项式法在第8章获得统一判据。}
{秩一不保证非零特征值存在，$\vu\perp\vv$就是反例。}
\section{相似与对角化判定}
\subsection{根的重数与特征向量数}
\begin{definition}{相似与相似对角化}
若$B=P\inv AP$且$P$可逆，称$A,B$相似。若可取$B=\diag(\lambda_i)$，称$A$可对角化。等式$AP=P\Lambda$表示$P$的每列为与$\Lambda$对角元同序的特征向量。
\end{definition}
\begin{theorem}{对角化的充要条件}
在特征多项式可分裂的数域上，$A$可对角化当且仅当有$n$个独立特征向量，当且仅当每个根均有$g_\lambda=a_\lambda$。不同特征值的特征向量线性无关；具有$n$个互异特征值必可对角化。相似保特征多项式、秩、迹、行列式，但这些量全部相同仍未必相似。
\end{theorem}
\begin{proof}
\label{der:5-2}
先证明不同特征值的向量无关。对$k$个互异根$\lambda_i$及非零向量$v_i$，以$k$归纳。$k=1$成立；若$\sum_i c_iv_i=0$，作用$A-\lambda_kI$得到
\[\sum_{i<k}c_i(\lambda_i-\lambda_k)v_i=0.\]
由归纳假设及非零差值，$c_1=\cdots=c_{k-1}=0$，原式再给$c_k=0$。同样的作用可应用于来自各特征子空间的向量，因此不同特征子空间之和为直和。

若有$n$个独立特征向量，拼成可逆$P$，按对应根排列$\Lambda$，则逐列$AP=P\Lambda$，左乘$P\inv$得对角化。反向若$P\inv AP=\Lambda$，$AP=P\Lambda$说明$P$的$n$列就是独立特征向量。分裂数域中$\sum a_\lambda=n$且$g_\lambda\le a_\lambda$，各特征空间直和维数为$\sum g_\lambda$，等于$n$恰当于每个$g_\lambda=a_\lambda$，得到所有等价判据。

若$B=P\inv AP$，则$\det(tI-B)=\det(P\inv(tI-A)P)=\det(tI-A)$，秩由可逆左右乘保持。迹满足$\tr(XY)=\sum_{ij}x_{ij}y_{ji}=\tr(YX)$，故$\tr B=\tr(APP\inv)=\tr A$；行列式取$t=0$或直接用乘积公式得到不变。相同的不变量不足以反推相似，例\ref{ex:5-8}给出具体反例。

\end{proof}
\begin{method}{重根用秩补足判断}
先求$\chi_A$，逐根算$n-\rank(A-\lambda I)$；每个几何重数等于代数重数才可对角化。构造$P$后检验$AP=P\Lambda$，比求$P\inv$再连乘省算。易错：重根可能可对角化，也可能不可。
\end{method}
\worked{构造对角化矩阵}{考研风格}{\star\star}{5-5}
{对角化例\ref{ex:5-1}中的$A$并求$A^m$，$m\ge0$。}
{把两个特征向量按根的顺序拼列。}
{\textbf{方法一：}以对应根2、3的向量拼$P=\begin{pmatrix}1&1\\0&1\end{pmatrix}$，逆为$\begin{pmatrix}1&-1\\0&1\end{pmatrix}$。逐列$AP=P\diag(2,3)$，故
\[A^m=P\diag(2^m,3^m)P\inv=\begin{pmatrix}2^m&3^m\\0&3^m\end{pmatrix}\begin{pmatrix}1&-1\\0&1\end{pmatrix}=\begin{pmatrix}2^m&3^m-2^m\\0&3^m\end{pmatrix}.\]
\textbf{方法二：}写$A^m=\begin{pmatrix}2^m&b_m\\0&3^m\end{pmatrix}$，$A^{m+1}=A^mA$给$b_{m+1}=2^m+3b_m,b_0=0$。设$b_m=3^m-2^m$，则$2^m+3(3^m-2^m)=3^{m+1}-2^{m+1}$，且初值成立，故归纳得到同式，包含$m=0$。}
{对角化同时处理所有幂；递推适合只追踪少量元素。}
{$P$列交换时$\Lambda$对角元必须同步交换。}
\worked{二重根的秩判定}{考研风格}{\star\star}{5-6}
{判断$A=\begin{pmatrix}2&1\\0&2\end{pmatrix}$是否可对角化。}
{唯一根2的代数重数是2。}
{\textbf{方法一：}$\chi_A(t)=(t-2)^2$，代数重数2。$A-2I=\begin{pmatrix}0&1\\0&0\end{pmatrix}$秩1，核方程$y=0$只有一个独立方向，几何重数1小于2，不能组成二阶特征基。
\textbf{方法二：}若存在$P\inv AP=\Lambda$且$\Lambda$对角，其两个对角元只能都是2，即$\Lambda=2I$。于是$A=P(2I)P\inv=2PP\inv=2I$，但题目$(1,2)$元为1，矛盾。因此不可对角化。}
{单特征值情形，标量矩阵反证尤其快。}
{有特征值不代表有足够特征向量。}
\worked{三阶重根判断}{考研风格}{\star\star}{5-7}
{判断$A=\begin{pmatrix}2&1&0\\0&2&0\\0&0&3\end{pmatrix}$是否可对角化。}
{只要一个特征值缺少特征向量，整体即失败。}
{上三角矩阵给$\chi_A(t)=(t-2)^2(t-3)$。根2时
$A-2I=\begin{pmatrix}0&1&0\\0&0&0\\0&0&1\end{pmatrix}$，其两条非零行独立，秩2，故$g_2=3-2=1<2$。根3时$A-3I=\begin{pmatrix}-1&1&0\\0&-1&0\\0&0&0\end{pmatrix}$秩2，$g_3=1$。全部特征空间直和维数$1+1=2<3$，所以不可对角化；缺少的一个方向不能从根3补齐。}
{秩判定比试图凑三个向量可靠。}
{三个对角元素并不自动给出三个独立特征向量。}
\worked{同谱仍不相似}{考研风格}{\star\star}{5-8}
{比较$D=\diag(1,1,2)$与$B=\begin{pmatrix}1&1&0\\0&1&0\\0&0&2\end{pmatrix}$是否相似。}
{寻找比迹和行列式更强的不变量。}
{\textbf{方法一：}$D-I=\diag(0,0,1)$秩1；$B-I=\begin{pmatrix}0&1&0\\0&0&0\\0&0&1\end{pmatrix}$秩2。若$B=P\inv DP$，则$B-I=P\inv(D-I)P$应同秩，矛盾。
\textbf{方法二：}$D$本身对角，可对角化；$B$在根1上的代数重数2，但$\ker(B-I)$仅$\Span(e_1)$，几何重数1，因此不可对角化。相似矩阵的特征空间由$P\inv$双射对应，不能出现这种差异。两矩阵虽然同迹4、同行列式2且同特征多项式，仍不相似。}
{秩族$\rank(A-\lambda I)$可识别部分Jordan结构。}
{二者特征多项式、迹与行列式全相同，不能据此认定相似。}
\section{实对称矩阵与正交对角化}
\subsection{谱定理及其几何意义}
\begin{theorem}{实对称谱定理}
实对称矩阵$A=A\T$的特征值全实；不同特征值的特征向量正交；存在实正交$Q$使
\begin{equation}Q\T AQ=\Lambda,\qquad A=\sum_i\lambda_i\vq_i\vq_i\T.\label{eq:spectral}\end{equation}
对非零$\vx$，Rayleigh商$R_A(\vx)=\frac{\vx\T A\vx}{\vx\T\vx}$满足$\lambda_{\min}\le R_A(\vx)\le\lambda_{\max}$。
\end{theorem}
\begin{proof}
\label{der:5-3}
先证明实对称$A$的根是实数。取复特征对$Az=\lambda z,z\ne0$，记$z^*$为共轭转置。$z^*z=\sum|z_i|^2>0$，且
\[\overline{z^*Az}=z^*A^*z=z^*Az\quad(A^*=A),\]
所以$z^*Az$为实数，$\lambda=\frac{z^*Az}{z^*z}$也实。若复向量$z=u+iv$对应这个实根，则$Au=\lambda u,Av=\lambda v$，$u,v$至少一个非零，所以可选实特征向量。

若$Au=\lambda u,Av=\mu v$，对称性给$\lambda u\T v=(Au)\T v=u\T Av=\mu u\T v$；当$\lambda\ne\mu$时得到正交。取一单位实特征向量$q_1$。对$x\perp q_1$，$q_1\T Ax=(Aq_1)\T x=\lambda_1q_1\T x=0$，所以$q_1^\perp$在$A$下不变。取其单位正交基，限制矩阵仍实对称，因为$B\T AB$的转置等于它本身。在$n-1$维限制上归纳得到单位正交特征基，补$q_1$即得全空间正交矩阵$Q$与$Q\T AQ=\Lambda$。$n=1$为归纳起点。

由$QQ\T=I$，$A=Q\Lambda Q\T=\sum_i\lambda_iq_iq_i\T$。令$x=Qy$，则
\[\dfrac{x\T Ax}{x\T x}=\dfrac{\sum_i\lambda_i y_i^2}{\sum_i y_i^2}.
\]
分母正，权重$\frac{y_i^2}{\sum_jy_j^2}$非负且和为1，因此商位于最小、最大特征值之间；取相应单位特征向量即可达到两端。等号向量恰落在对应极值特征子空间中。

\end{proof}
\begin{method}{重根先在特征空间内正交化}
不同特征值之间天然正交；同一重根下任选基，作Schmidt后单位化；拼成$Q$。检查$Q\T Q=I$与$AQ=Q\Lambda$。不能对一般非对称矩阵强行要求正交对角化。
\end{method}
\worked{二阶正交对角化}{基础例题}{\star\star}{5-9}
{正交对角化$A=\begin{pmatrix}2&1\\1&2\end{pmatrix}$。}
{对称的和、差方向容易观察。}
{\textbf{方法一：}$A(1,1)\T=(3,3)\T$、$A(1,-1)\T=(1,-1)\T$，两向量正交且长度$\sqrt2$。取$q_1=\frac{(1,1)\T}{\sqrt2},q_2=\frac{(1,-1)\T}{\sqrt2}$，拼$Q$后$Q\T Q=I$，由$Aq_1=3q_1,Aq_2=q_2$得$Q\T AQ=\diag(3,1)$。
\textbf{方法二：}$\chi_A(t)=(t-2)^2-1=t^2-4t+3=(t-3)(t-1)$；根3的核方程$-x+y=0$、根1的核方程$x+y=0$，分别给同一组和差向量，再单位化。两法都提供实际正交矩阵$Q=\frac1{\sqrt2}\begin{pmatrix}1&1\\1&-1\end{pmatrix}$。}
{观察法适合对称结构，特征方程法更一般。}
{只用未单位化向量拼列得的是相似矩阵，通常不是正交矩阵。}
\worked{重根特征空间的正交基}{考研风格}{\star\star}{5-10}
{正交对角化$A=I_3+\bm1\bm1\T$。}
{全1方向与分量和零平面分开。}
{\textbf{方法一：}令$q_1=\frac{\bm1}{\sqrt3}$，因为$\bm1\T\bm1=3$，$Aq_1=4q_1$。取$q_2=\frac{(1,-1,0)\T}{\sqrt2},q_3=\frac{(1,1,-2)\T}{\sqrt6}$，二者分量和零，因此$Aq_2=q_2,Aq_3=q_3$。三个点积均零，长度均1，所以$Q=(q_1,q_2,q_3)$正交，$Q\T AQ=\diag(4,1,1)$。
\textbf{方法二：}解$(A-I)x=0$得到$x_1+x_2+x_3=0$，选基$u=(1,-1,0),v=(1,0,-1)$。$u\T v=1,\norm u^2=2$，减投影$v-\frac{u}{2}=(\frac{1}{2},\frac{1}{2},-1)$，放大2并单位化即$q_3$；$u$单位化即$q_2$。再补$q_1$，得到同一$Q$，根1的两个方向均已列全。}
{特征空间内的基不唯一，正交对角化矩阵也不唯一。}
{根1重数为2，不能只提供一个对应向量。}
\worked{不同根正交的直接证明}{考研风格}{\star\star}{5-11}
{设实对称$A$有特征对$(\lambda,\vu),(\mu,\vv)$，$\lambda\ne\mu$，证明正交。}
{用对称性把$A$移到内积另一边。}
{从$Au=\lambda u$出发，将内积中的$A$移到另一边：
\[\lambda u\T v=(Au)\T v=u\T A\T v=u\T Av=\mu u\T v.\]
中间$A\T=A$是关键。相减得$(\lambda-\mu)u\T v=0$，而$\lambda\ne\mu$，可除以差值，得到$u\T v=0$。若根相同，差值为零，方程不提供正交信息，所以同一特征空间内还须主动选正交基。}
{对称性是论证的关键，普通相似矩阵不能照搬。}
{重根时这条等式变成$0=0$，不保证任取两特征向量正交。}
\worked{Rayleigh商的两法极值}{竞赛风格}{\star\star}{5-12}
{求$x^2+y^2=1$时$2x^2+2xy+2y^2$的最值。}
{这是例\ref{ex:5-9}矩阵的二次型。}
{\textbf{方法一：}矩阵$A=\begin{pmatrix}2&1\\1&2\end{pmatrix}$在单位和差基下变成$\diag(3,1)$。令$(x,y)\T=Q(u,v)\T$，约束变成$u^2+v^2=1$，目标$3u^2+v^2=1+2u^2$，范围$[1,3]$。最大要求$v=0$，对应$x=y=\pm\frac{1}{\sqrt2}$；最小要求$u=0$，对应$x=-y=\pm\frac{1}{\sqrt2}$。
\textbf{方法二：}目标为$2(x^2+y^2)+2xy=2+2xy$。由$(x-y)^2\ge0,(x+y)^2\ge0$分别得$2xy\le1,2xy\ge-1$，同样给范围和等号条件。极值均实际取得。}
{谱法推广高维；二元不等式在小题更快。}
{极值必须同时给出满足约束的等号向量。}
\begin{insight}{特征向量是变换后仍留在原直线上的方向}
特征向量的长度可变、箭头可反向，但它张成的直线保持不变。一般向量则会混合多个特征方向，方向随之改变。对$A=\begin{pmatrix}2&1\\1&2\end{pmatrix}$，和方向被放大3倍，差方向不变，而$e_1$被送到$(2,1)\T$，并不是它的倍数。
\end{insight}
\begin{figure}[H]\centering
\begin{tikzpicture}[scale=1.1]
\draw[laaxis](-1.3,0)--(3.7,0)node[right]{$x$};\draw[laaxis](0,-1.6)--(0,3.6)node[above]{$y$};
\draw[dashed,Teal!50](-1,-1)--(3.1,3.1);\draw[dashed,Gold!70](-1.2,1.2)--(1.2,-1.2);
\draw[lavec,Teal!40,line width=2.3pt](0,0)--(3,3)node[above]{$A\vu=3\vu$};
\draw[lavec,Teal](0,0)--(1,1)node[above left]{$\vu$};
\draw[lavec,Gold](0,0)--(1,-1)node[below right]{$A\vv=\vv$};
\draw[lavec,Ink](0,0)--(1,0)node[below]{$e_1$};
\draw[lavec,Rust](0,0)--(2,1)node[right]{$Ae_1$};
\node[lanote,anchor=west]at(4,.7){沿特征方向：只缩放\\[3pt]一般方向：发生混合};
\end{tikzpicture}
\caption{和差方向是不变直线；并非每个向量都是特征向量。}\label{fig:eigen-directions}
\end{figure}
\begin{clevertrick}{用两个投影直接求所有幂}
令$E_+=\frac12\begin{pmatrix}1&1\\1&1\end{pmatrix}$、$E_-=I-E_+=\frac12\begin{pmatrix}1&-1\\-1&1\end{pmatrix}$。对$x=(s,t)\T$，两投影分别给$E_+x=\frac{(s+t)(1,1)\T}{2}$、$E_-x=\frac{(s-t)(1,-1)\T}{2}$，确实是和、差分量。直接乘法给$E_+^2=E_+$；由$E_-=I-E_+$得
\[E_-^2=I-2E_++E_+^2=E_-,\qquad E_+E_-=E_-E_+=E_+-E_+^2=0.\]
又$3E_++E_-=I+2E_+=\begin{pmatrix}2&1\\1&2\end{pmatrix}=A$。对$m\ge1$，乘方展开中凡同时含两种投影的项都为零，纯投影项由幂等性化简，因此
\[A^m=3^mE_++E_-=\dfrac12\begin{pmatrix}3^m+1&3^m-1\\3^m-1&3^m+1\end{pmatrix}.\]
$m=0$时右端为$E_++E_-=I$，仍成立；负整数也可先核验$A\inv=\frac{E_+}{3}+E_-$再同样推导。这把向量分成两个互不干扰的分量，分别缩放后相加，省去显式求$P\inv$。
\textbf{易错提示：}负特征值会反向；零特征值把整条特征方向压到零，并不意味着非零向量失去作为特征向量的资格。
\end{clevertrick}


\section{Jordan结构与谱方法边界}
\subsection{选学提高：广义特征向量}
\begin{extension}
本节为选学提高。考研核心只需熟练判断是否有足够特征向量；Jordan块、广义链与核维数将在第\ref{ch:advanced}章进一步联系最小多项式。
\end{extension}
\begin{definition}{Jordan块与广义链}
$J_k(\lambda)=\lambda I_k+N_k$，$N_k$的第一超对角元为1，其余为0。若$\vv_1\ne0$，$(A-\lambda I)\vv_1=0$，且$(A-\lambda I)\vv_{j+1}=\vv_j$对$j=1,\ldots,k-1$成立，称$\vv_1,\ldots,\vv_k$为长度$k$的广义特征链。一般称满足$(A-\lambda I)^s v=0$的非零向量$v$为广义特征向量，其中$s$为某个正整数。

链向量必无关：若存在非零关系，取最大下标$j$使其系数$c_j\ne0$，对该关系作用$(A-\lambda I)^{j-1}$，前面各项全为零，只剩$c_j\vv_1=0$，与$\vv_1\ne0$矛盾。因此链张成$k$维不变子空间；按链序作基，该限制变换的矩阵才是$J_k(\lambda)$。
\end{definition}
\begin{theorem}{Jordan存在与幂公式（标准结论）}
若特征多项式在数域中分裂，则$A$相似于Jordan块的直和，块大小在重排意义下唯一。对非负整数$m$，
\begin{equation}J_k(\lambda)^m=\sum_{j=0}^{\min(m,k-1)}\binom mj\lambda^{m-j}N_k^j.\label{eq:jpower}\end{equation}
当$\lambda=0$时，直接读取为$J_k(0)^m=N_k^m$；特别$m=0$时为$I$，不另计算未约定的标量$0^0$。以下从广义特征空间的分离与幂零链的构造证明存在性和唯一性。
\end{theorem}
\begin{proof}
\label{der:jordan-full}
\textbf{一、把不同特征值分离。}在分裂数域中写$\chi_A(t)=\prod_{i=1}^r f_i(t)$，$f_i(t)=(t-\lambda_i)^{a_i}$，根$\lambda_i$互异，令$F_i=\frac{\chi_A}{f_i}$。$f_i,F_i$互素，欧几里得算法得到多项式$u_i,v_i$使$u_iF_i+v_if_i=1$。其依据是反复带余除法求出最大公因式1，再逆序回代表示为两个原多项式的组合。

令$E_i=u_i(A)F_i(A)$。对每个$j$，多项式$\sum_i u_iF_i-1$模$f_j$的余数为零：当$i\ne j$，$F_i$含$f_j$；当$i=j$，由Bezout等式余数为1。两两互素的$f_j$之积因此整除它。凯莱--哈密顿定理（证明见推导\ref{der:8-2}）给
\[\sum_iE_i=I,\qquad f_i(A)E_i=u_i(A)\chi_A(A)=0.\]
于是每个$x$均为$\sum_i E_ix$，其中$E_ix\in V_i:=\ker f_i(A)$。在$V_i$上，$E_i=I$；在$V_j,j\ne i$上，$E_i=0$。若$\sum_i x_i=0,x_i\in V_i$，作用$E_j$得$x_j=0$，所以$V=\bigoplus_iV_i$。限制$N_i=(A-\lambda_iI)|_{V_i}$满足$N_i^{a_i}=0$。

\textbf{二、为任意幂零映射构造链基。}设$N^s=0$且$s$最小，令$K_k=\ker N^k,d_k=\dim K_k$，其中$K_0=0,K_s=K_{s+1}=V$。对$k=1,\ldots,s$选补空间$H_k$使
\[K_k=(K_{k-1}+NK_{k+1})\oplus H_k.\]
这是在有限维空间中把子空间基扩充成$K_k$的基。对$H_k$的每个基向量$h$，列出
\[N^{k-1}h,\ N^{k-2}h,\ \ldots,\ h.\]
因为$h\notin K_{k-1}$，首向量非零；又$N^kh=0$，所以这确为长度$k$的链。

记全部链向量的生成空间为$S$，则$NS\subseteq S$。由$NK_{s+1}=NV\subseteq K_{s-1}$得$K_s\subseteq S+K_{s-1}$。若已经有$K_{k+1}\subseteq S+K_k$，则
\[K_k=K_{k-1}+NK_{k+1}+H_k\subseteq K_{k-1}+NS+NK_k+S\subseteq K_{k-1}+S.\]
从$s$向下归纳，再逐级代入，得到$V=K_s\subseteq S$，所以链向量生成全空间。

还须证明没有冗余。因为$NK_{k+1}\cap K_{k-1}=NK_k$：若$Nv\in K_{k-1},v\in K_{k+1}$，则$N^kv=0$，所以$v\in K_k$；反向包含立即成立。又$\dim NK_j=d_j-d_1$，故
\[\begin{aligned}
\dim H_k&=d_k-\dim(K_{k-1}+NK_{k+1})\\
&=d_k-\{d_{k-1}+(d_{k+1}-d_1)-(d_k-d_1)\}\\
&=2d_k-d_{k-1}-d_{k+1}.
\end{aligned}\]
令$\delta_k=d_k-d_{k-1},\delta_{s+1}=0$，全部链向量数为
\[\sum_{k=1}^s k\dim H_k=\sum_{k=1}^s k(\delta_k-\delta_{k+1})=\sum_{k=1}^s\delta_k=d_s=\dim V.\]
已生成全空间且向量数等于维数，故它们无关，构成链基。

\textbf{三、得到Jordan形并证明唯一。}在链序$N^{k-1}h,\ldots,h$下，$N$的矩阵是超对角为1的$J_k(0)$，加回$\lambda_i I$便是$J_k(\lambda_i)$。不同广义特征空间的链基拼接成全空间基，因此得到Jordan标准形。任一长度$k$的块对$\dim\ker N^j$贡献$\min(j,k)$，其相邻差贡献1恰在$j\le k$时成立。所以长度至少$j$的块数被$d_j-d_{j-1}$唯一决定，恰为$j$的块数再取差得$2d_j-d_{j-1}-d_{j+1}$。这些维数由原映射决定，故块大小仅能重排，不能改变。

\textbf{四、幂公式。}$\lambda I$与$N_k$可交换，有限二项式展开给
\[J_k(\lambda)^m=\sum_{j=0}^m\binom mj\lambda^{m-j}N_k^j.
\]
因$N_k^j=0$对$j\ge k$，和截断到$\min(m,k-1)$。$m=0$时矩阵零次幂为$I$。$\lambda=0$时原矩阵就是$N_k$，故幂直接为$N_k^m$，按这个结果读取显示公式，避免另行计算标量$0^0$。

\end{proof}
\begin{method}{缺少向量时改求链}
先求核中的非零$\vv_1$，再解$(A-\lambda I)\vv_2=\vv_1$。不是每个非零核向量都能延成长链，需检查它是否在像中。Jordan存在要求特征多项式在所用数域分裂；如例\ref{ex:5-14}的直角旋转矩阵，在实数域就没有特征向量。一般二维旋转中，转角为$0$或$\pi$时分别为$I$或$-I$，须与该例区别。
\end{method}
\worked{二阶Jordan块的幂}{竞赛风格}{\star\star\star}{5-13}
{求$A=\begin{pmatrix}2&1\\0&2\end{pmatrix}$的$m$次幂，$m\ge1$。}
{它等于$2I+N$且$N^2=0$。}
{\textbf{方法一：}写$A=2I+N$，$N=\begin{pmatrix}0&1\\0&0\end{pmatrix}$且$N^2=0$。可交换二项式展开只保留$j=0,1$：$A^m=2^mI+m2^{m-1}N$，故为$\begin{pmatrix}2^m&m2^{m-1}\\0&2^m\end{pmatrix}$。
\textbf{方法二：}令右上元$b_m$，$A^{m+1}=A^mA$给$b_{m+1}=2^m+2b_m$，$b_1=1$。除以$2^{m+1}$，令$c_m=\frac{b_m}{2^m}$，得到$c_{m+1}=\frac{1}{2}+c_m,c_1=\frac{1}{2}$，故$c_m=\frac{m}{2}$，即$b_m=m2^{m-1}$。此题$m\ge1$，不涉及负幂。}
{不可对角化仍可高效求幂；额外的$m$来自Jordan链。}
{不能只把对角元取幂而丢失右上元。}
\worked{数域决定可对角化性}{考研风格}{\star\star}{5-14}
{研究$R=\begin{pmatrix}0&-1\\1&0\end{pmatrix}$在实数域和复数域的对角化。}
{特征方程为$t^2+1=0$。}
{$\chi_R(t)=\det\begin{pmatrix}t&1\\-1&t\end{pmatrix}=t^2+1$。实数域没有根，因此没有实特征向量，不能实相似对角化，更不能实正交对角化。复数域根$i,-i$互异：$R(1,-i)\T=(i,1)\T=i(1,-i)\T$，$R(1,i)\T=(-i,1)\T=-i(1,i)\T$。取$P=\begin{pmatrix}1&1\\-i&i\end{pmatrix}$，$\det P=2i\ne0$，所以$P\inv RP=\diag(i,-i)$。数域改变了可用的基，而不是改变原矩阵元素。}
{“正交矩阵”不等于“实对称矩阵”，只有后者保证实正交对角化。}
{实数域没有特征值，复数域却有完整特征基；先说明数域，再谈对角化。}
\worked{乘积的谱不能随意相乘}{竞赛风格}{\star\star}{5-15}
{已知$A$谱为$1,2$，$B$谱为$3,4$，能否唯一确定$AB$的谱？}
{改变$B$特征方向的配对方式。}
{固定$A=\diag(1,2)$，选$B_1=\diag(3,4)$与$B_2=\diag(4,3)$。二者谱作为无序多重集均为$\{3,4\}$，但
\[AB_1=\diag(3,8),\qquad AB_2=\diag(4,6).\]
乘积的两个谱集合不同，故已知条件不足以唯一确定。两组中矩阵甚至都可交换，仍须知道哪个$B$特征值与哪个$A$特征方向配对。可确定的只是乘积行列式$(1\cdot2)(3\cdot4)=24$；迹分别11与10，确实不同。}
{共同特征基下才可按相应向量配对相乘；仅知两个无序谱不够。}
{可确定的是$\det(AB)=\det A\det B=24$，但迹仍可不同。}
\worked{由核维数读Jordan块}{竞赛风格}{\star\star\star}{5-16}
{设$N$是四阶幂零矩阵，$\dim\ker N=2,\dim\ker N^2=3,\dim\ker N^3=4$，求Jordan块大小。}
{核维数每次增量等于长度至少为该次幂的链数。}
{$d_0=0,d_1=2,d_2=3,d_3=4$，因空间仅四维，$d_4=4$。相邻增长量为$2,1,1,0$，分别是大小至少1、2、3、4的块数。因此大小恰1的块数$2-1=1$，恰2为$1-1=0$，恰3为$1-0=1$。所以$N\sim J_3(0)\oplus J_1(0)$。检查总阶$3+1=4$，核维数依次$1+1=2,2+1=3,3+1=4$，与数据吻合，故结构唯一。}
{秩或核维数比特征多项式提供更细的结构信息。}
{仅知$\chi_N=t^4$，可能的Jordan块不唯一。}
\begin{insight}{同一个特征值也可能藏着剪切}
矩阵$J=\begin{pmatrix}1&1\\0&1\end{pmatrix}$所有特征值都是1，但它并非恒等变换。它只在水平轴上保持向量不变；对$e_2$连续作用，每一步都沿水平轴多走一格。特征多项式没有记录下这条“传递链”，Jordan结构补上了它。
\end{insight}
\begin{figure}[H]\centering
\begin{tikzpicture}[scale=1.25]
\draw[laaxis](-.3,0)--(4.4,0)node[right]{$x$};\draw[laaxis](0,-.2)--(0,1.8)node[above]{$y$};
\draw[dashed,Teal!50](-.1,1)--(3.6,1);
\foreach \m in {0,1,2,3}{\draw[lavec,Teal!75](0,0)--(\m,1);\node[lapt]at(\m,1){};\node[font=\small,above]at(\m,1.08){$m=\m$};}
\draw[-{Stealth},Gold,thick](.2,1.45)--(2.8,1.45)node[midway,above]{每次增加$e_1$};
\node[lanote]at(5,1){$J^m e_2=(m,1)\T$\\[3pt]$E_1=\Span(e_1)$};
\end{tikzpicture}
\caption{特征值相同仍可产生剪切；一次项$m$反映长度2的Jordan链。}\label{fig:jordan-shear}
\end{figure}
\begin{clevertrick}{唯一特征值时，先看是否为标量阵}
在特征多项式分裂的数域中，若$n$阶矩阵只有特征值$\lambda$，它可对角化当且仅当$A=\lambda I$。必要性：若$A=PDP\inv$且$D$对角，所有对角元均只能是$\lambda$，故
\[A=P(\lambda I)P\inv=\lambda PP\inv=\lambda I.\]
充分性：若$A=\lambda I$，任意一组基的每个向量都被乘$\lambda$，其表示本来就是对角矩阵。例\ref{ex:5-6}只有根2而非$2I$，所以立即排除对角化。
图中的$J=I+N$满足$N^2=0,Ne_2=e_1$，因此$J^me_2=e_2+m e_1=(m,1)\T$，长度$\sqrt{m^2+1}$随$m\to\infty$无界。这个显式计算说明“全部特征值模不超过1”本身不足以保证矩阵幂有界。
\textbf{易错提示：}图中所有特征值的模均为1，并不能推出幂保持有界；$J^m e_2$的长度仍随$m$增大。谱相同也不能推出动力学行为相同。
\end{clevertrick}


\begin{warning}
相似要求一套基对定义域和值域同时改变，不是任意左右等价变换。相同谱不是充分相似条件。不同特征值向量必独立，对称条件才给正交。实对称矩阵的秩等于非零特征值按代数重数计数的总数；普通矩阵的秩不能这样数，例如非零幂零矩阵。
\end{warning}
\section{速算卡：伴随矩阵的特征值与迹}\label{r2:adjugate}
本书 $A^*$ 表示古典伴随矩阵（adjugate），不是共轭转置。设 $A\in\mathbb C^{n\times n}$，特征值 $\lambda_1,\ldots,\lambda_n$ 均按代数重数计。

\begin{theorem}{可逆时的倒数谱与任意方阵的乘积谱}
若 $A$ 可逆，则 $A^*=(\det A)A^{-1}$，因此 $A^*$ 的特征值为
\[
 \frac{\det A}{\lambda_i}=\prod_{j\ne i}\lambda_j\quad(i=1,\ldots,n).
\]
并且对任意 $A$，包括奇异矩阵，均有
\[
 \operatorname{tr}(A^*)=\sum_{i=1}^n\prod_{j\ne i}\lambda_j.
\]
奇异时只使用右侧不含除法的表达，不能把零特征值代入 $\det A/\lambda_i$。事实上伴随矩阵的谱仍是这 $n$ 个乘积，按重数计；这一结论不要求 $A$ 可对角化。
\end{theorem}
\begin{proof}
可逆情形由 $AA^*=(\det A)I$ 直接得到。任意 $A$ 可在复数域相似上三角化；伴随变换与相似相容，可先对 $A+tI$ 的可逆取值证明，再利用伴随各元素是原元素的多项式延伸到 $t=0$。上三角矩阵的伴随仍上三角，其第 $i$ 个对角元就是其余对角元之积。因此得到谱与迹公式。也可从特征多项式的 $t$ 一次项看出迹公式。
\end{proof}

\begin{example}{可逆与奇异两种速算}{考研训练}{r2-adjugate}
若三阶 $A$ 的特征值为 $1,2,3$，则 $\det A=6$，伴随特征值为 $6,3,2$，迹为 $11$。若改为 $0,2,3$，伴随谱为 $6,0,0$，迹为 $6$；虽 $\det A=0$，伴随未必为零。
另一例 $A=\begin{pmatrix}2&1\\0&2\end{pmatrix}$ 不可对角化，但 $A^*=\begin{pmatrix}2&-1\\0&2\end{pmatrix}$ 的谱仍为 $2,2$，迹为四。
\end{example}

\paragraph{与秩判据一起使用}
当 $n\ge2$ 时，$\operatorname{rank}A=n,n-1,\le n-2$ 分别对应 $\operatorname{rank}A^*=n,1,0$。但秩为一的伴随也可能所有特征值为零，不能把“秩一”当作“有一个非零特征值”。例如非零二阶幂零阵 $A=\begin{pmatrix}0&1\\0&0\end{pmatrix}$ 的伴随为 $-A$，秩一而谱全零。
自测：四阶 $A$ 的谱为 $1,1,2,4$。答：$\det A=8$，伴随谱为 $8,8,4,2$，迹 $22$。$n=1$ 时按空乘积为一的约定，$A^*=[1]$，不要把 $n\ge2$ 的秩分类直接搬用。


\section*{章末分层习题}\addcontentsline{toc}{section}{章末分层习题}
\exercise{5-1}{基础}{求$\begin{pmatrix}1&2\\0&3\end{pmatrix}$的特征值与特征空间。}
\exercise{5-2}{基础}{若三阶矩阵特征值为$-1,2,2$，求迹、行列式与$\tr(A^2)$。}
\exercise{5-3}{基础}{判断$\diag(2,2,3)$是否可对角化。}
\exercise{5-4}{考研风格}{判断$\begin{pmatrix}1&t\\0&1\end{pmatrix}$何时可对角化。}
\exercise{5-5}{考研风格}{正交对角化$\begin{pmatrix}3&1\\1&3\end{pmatrix}$。}
\exercise{5-6}{考研风格}{若$n$阶$A\sim B$，证明对任意$\lambda\in\F$及非负整数$k$，$\rank\bigl((A-\lambda I)^k\bigr)=\rank\bigl((B-\lambda I)^k\bigr)$。}
\exercise{5-7}{竞赛拓展}{若$A^2=A$，证明可对角化且谱只可能为0、1。}
\exercise{5-8}{竞赛拓展}{对非负整数$m$，求$J_3(1)^m$的各上三角元素，并指出$m=0,1$时哪些元素为零。}
\exercise{5-9}{竞赛拓展}{实对称$A$满足所有特征值位于$[2,5]$，证明$2\norm{\vx}^2\le\vx\T A\vx\le5\norm{\vx}^2$。}
\section*{章末方法速查}\addcontentsline{toc}{section}{章末方法速查}
\begin{quicktable}
特征方程&方阵&未知谱&统一$tI-A$符号\\
迹与行列式&谱计算核验&已知部分根&按代数重数计\\
秩判重根&特征多项式分裂&重根出现&算每个几何重数\\
拼特征基&向量数足够&构造相似矩阵&列与根同序\\
正交对角化&实对称&重根空间&先正交再单位化\\
Jordan链&选学，分裂数域&缺少特征向量&检查链可延长性\\
\end{quicktable}
